\scanpage{9}\mspage{- IV-\handwritten{1131} -}
% Single long, near-vertical pencil stroke through the setting copy.
\margin{\handwritten{\unsure{?}}}
existe , en vertu de (21.9.12) un $\underline{\mathrm O}_X$-Module inversible $\underline L$ dont \deleted{\illegible{1}}
l'image réciproque dans $X'_0$ est isomorphe à $\underline L_0$ , donc est ample . Puisque tout voisinage de $s_0$ dans $S$ est nécessairement $S$ tout entier ,on déduit alors de (9.6.4) que $\underline L$ est un $\underline{\mathrm O}_X$-Module ample , d'où la conclusion \deleted{\illegible{1}} (\uline{II},5.3.1 et \uline{II},5.5.3)\deleted{\illegible{1}} .

\uline{21.10 . Factorialité des anneaux locaux réguliers} .

\uline{Théorème} (21.10.1) \inserted{(Auslander-Buchsbaum)} .- \uline{Un anneau local noethérien régulier est factorie\handwritten{l}}.
% slanted
\margin{\handwritten{Créditer Kaplansky de la \unsure{démonstration|démonstr.} donnée ici.}}

Soit $A$ un anneau local noethérien régulier de dimension $n$ ; nous allon\handwritten{s} raisonner par récurrence sur $n$ . Pour $n=\unsure{0}$,$A$ est un corps et pour $n=1$ ,
un anneau de valuation discrète , donc principal (et \uline{a fortiori} factoriel) . Supposons donc $n\handwritten{\geqslant}2$ et le théorème démontré pour des anneaux réguliers d\handwritten{e} dimension $\handwritten{<}n$ . Posons $X=\operatorname{Spec}(A)$ , désignons par $a$ le point fermé de $X$ et posons $U=X-\{a\}$ . En tout point $y\in U$ , on a \deleted{\illegible{1}},
$\dim(\underline{\mathrm O}_{X,y})\handwritten{\leqslant}n-1$ , et puisque $A$ est régulier , les anneaux $\underline{\mathrm O}_{X,y}$ le sont aussi (\uline{0},17.3.2), donc l'hypothèse de récurrence entraîne qu'ils sont factoriels . De plus on a $\operatorname{prof}(A)=\dim(A)\handwritten{\geqslant}2$ puisque $A$ est régulier , donc de Cohen-Macaulay (\uline{0},17.1.3) . Utilisant (21.6.14) , on est ramen\handwritten{é} à prouver que $\operatorname{Pic}(U)=0$ . Considérons donc un $\underline{\mathrm O}_U$-Module inversible $\underline L$ , et prouvons qu'il est isomorphe à $\underline{\mathrm O}_U$ . \deleted{Puisque $\dim(A)\handwritten{\geqslant}2$ et que $A$ est régulier ,} \handwritten{I}l résulte de \deleted{(5.11.6)} \inserted{\handwritten{I 9.4.5.}} qu'il existe un $\underline{\mathrm O}_X$-Module \uline{cohérent} $\underline F$ tel que $\underline F\handwritten{\mid}U=\underline L$ . Comme $A$ est régulier , donc de dimension cohomologique finie (\uline{0},17.3.1) , il existe une réso\handwritten{l}ution gauche \uline{finie} de $F$ :
\[
0\handwritten{\longleftarrow}\underline F\handwritten{\longleftarrow}\underline{\mathrm O}_X^{n_1}\handwritten{\longleftarrow}\underline{\mathrm O}_X^{n_2}\handwritten{\longleftarrow}\mathinner{\ldotp\ldotp\ldotp\ldotp}\handwritten{\longleftarrow}\underline{\mathrm O}_X^{n_h}\handwritten{\longleftarrow}0
\]
(\uline{0},17.2.8 et \uline{0},17.2.2,(iii)) . Par restriction à $U$ , on a donc une résolution finie
\begingroup\makeatletter\tagsleft@true\makeatother
\[
0\handwritten{\longleftarrow}\underline L\handwritten{\longleftarrow}\underline{\mathrm O}_U^{n_1}\handwritten{\longleftarrow}\underline{\mathrm O}_U^{n_2}\handwritten{\longleftarrow}\mathinner{\ldotp\ldotp\ldotp\ldotp}\handwritten{\longleftarrow}\underline{\mathrm O}_U^{n_h}\handwritten{\longleftarrow}0
\tag*{\handwritten{(21.10.1.1)}}
\]
\endgroup

Le théorème résultera alors des considérations générales suivantes \handwritten{\unsure{;|:}}

Sur un espace \deleted{\illegible{1}} annelé $X$ , soit $\underline E$ un $\underline{\mathrm O}_X$ -Module localement libre
% archive sheet number: 8
